\documentclass[11pt,a4paper]{article}

\usepackage[margin=2.5cm]{geometry}
\usepackage{amsmath,amssymb}
\usepackage{graphicx}
\usepackage{booktabs}
\usepackage{hyperref}
\usepackage{xcolor}
\usepackage{enumitem}
\usepackage{microtype}
\usepackage{parskip}
\usepackage{titlesec}
\usepackage{caption}
\usepackage{subcaption}
\usepackage{multirow}
\usepackage{array}
\usepackage{longtable}
\usepackage{lmodern}
\usepackage{csquotes}
\usepackage[round]{natbib}

\hypersetup{
  colorlinks=true,
  linkcolor=blue!60!black,
  citecolor=green!50!black,
  urlcolor=blue!60!black,
}

\titleformat{\section}{\large\bfseries}{\thesection}{1em}{}
\titleformat{\subsection}{\normalsize\bfseries}{\thesubsection}{1em}{}
\titleformat{\subsubsection}{\normalsize\itshape}{\thesubsubsection}{1em}{}

\newcommand{\margintext}[1]{\textcolor{gray}{\small #1}}
\newcommand{\todo}[1]{\textcolor{red}{\textbf{TODO:} #1}}
\newcommand{\layerz}{$\ell{=}0$}
\newcommand{\layerel}{$\ell{=}11$}

% ----------------------------------------------------------------
\title{\textbf{Progress Report: Cross-Modal Information Flow in\\
Multimodal Large Language Models}\\[0.5em]
\large Attention Knockout Reproduction, SAE Feature Ablation,\\
and Methodological Analysis of a Persistent Null Result}

\author{Ron Kremer}
\date{April 2026}
% ----------------------------------------------------------------

\begin{document}
\maketitle
\tableofcontents
\newpage

% ================================================================
\section{Introduction and Research Motivation}
\label{sec:intro}
% ================================================================

Multimodal large language models (MLLMs) such as LLaVA~\citep{liu2023llava} interleave visual tokens produced
by a vision encoder with textual token representations inside a shared transformer decoder.
Despite strong empirical performance on visual question answering (VQA), the precise computational
mechanisms by which visual information is transferred into language representations remain poorly
understood.
A central open question is: \emph{at which transformer layers, and through which attention
pathways, does visual content causally influence the model's textual predictions?}

The prior work ``Cross-Modal Information Flow in Multimodal Large Language
Models''~\citep{crossmodal2024} introduced a causal probing technique—\emph{attention knockout}—to
answer this question, demonstrating that a small number of early-to-middle transformer layers carry
the dominant fraction of image-to-text information transfer in LLaVA.
The present project extends that work in two directions:
(1)~reproducing and extending the attention knockout findings to a larger dataset and an additional
attention-flow direction; and
(2)~training Sparse Autoencoders (SAEs) at the identified layers to decompose the causal information
flow into individual interpretable features, testing whether specific feature subsets are causally
necessary for visual attribute binding.

This report documents the full experimental trajectory from initial reproduction through 18 SAE
ablation experiments, a diagnostic ceiling analysis, and a series of methodological fixes.
We report not only the positive results from attention knockout but also a comprehensive
null-result narrative from the SAE ablation phase, including a principled analysis of why the null
arose and the interventions we have applied or planned.

\subsection{Model and Task}

All experiments use \textbf{LLaVA-v1.5-7b} (\texttt{liuhaotian/llava-v1.5-7b}), a
7-billion-parameter MLLM built on Vicuna-7B (LLaMA-2) with a CLIP ViT-L/14@336 vision encoder and a
two-layer MLP visual projection head. The transformer decoder has 32 layers with hidden dimension
$d_{\mathrm{model}} = 4096$ and 32 attention heads. Image inputs are tokenised into 576 visual
patch tokens (via the \texttt{prepare\_inputs\_labels\_for\_multimodal} method) which are inserted
at the position of the special \texttt{IMAGE\_TOKEN\_INDEX = -200} placeholder.

The evaluation task is \textbf{ChooseAttr}, a forced-choice VQA benchmark derived from the GQA
validation set. Each item consists of an image, an object, and two candidate attribute values (e.g.,
\emph{``Is the car red or blue?''}) where exactly one option is correct. Performance is measured as
the log-probability \emph{margin}: $\Delta = \log P(\text{true option}) - \log P(\text{false
option})$, normalised across the full label string. The dataset comprises 1\,000 items spanning five
attribute categories (color, material, shape, size, state); as detailed in
Section~\ref{sec:dataset}, the \textbf{color} category (602 items) is the primary evaluation subset
due to its size and homogeneity.


% ================================================================
\section{Dataset Characterisation}
\label{sec:dataset}
% ================================================================

Prior to running experiments we conducted a thorough quality audit of the ChooseAttr dataset.

\subsection{Global Properties}

The full dataset contains 1\,000 rows drawn from GQA validation images. We verified that:
\begin{itemize}[noitemsep]
  \item There are no duplicate \texttt{question\_id} entries.
  \item The \texttt{answer} column equals the true option for all 1\,000 rows.
  \item Option order is balanced (501 true-first, 499 false-first), ruling out position bias.
  \item The false option is never drawn from the object's own attribute list (verified for all 602
    color rows); foils are plausible but incorrect distractors.
  \item 49 rows cover non-attribute tasks (pose, activity, face expression) and are correctly
    excluded from attribute-binding analyses.
\end{itemize}

\subsection{Per-Category Characterisation}

Table~\ref{tab:dataset} summarises the per-category breakdown.

\begin{table}[ht]
\centering
\caption{ChooseAttr dataset composition by attribute category.}
\label{tab:dataset}
\begin{tabular}{lrrll}
\toprule
Category & Rows & Runnable? & GQA attribute types & Notes \\
\midrule
color    & 602  & \checkmark & color (homogeneous) & $\sim$58\% involve black/white options; \\
         &      &            &                     & 20.6\% tiny objects ($<$1\% image area) \\
material &  96  & $\times$   & material            & clean but below threshold \\
shape    &  15  & $\times$   & shape               & critically small; 9/15 from three pairs \\
size     &  87  & $\times$   & size, length, height, & heterogeneous dimension bundle \\
         &      &            & depth, weight, \ldots & \\
state    & 151  & \checkmark & weather, cleanliness, & 77\% are generic ``choose X|Y'' items \\
         &      &            & + 116 generic items  & with no GQA attribute type \\
\bottomrule
\end{tabular}
\end{table}

\noindent The analysis also revealed a \textbf{language-side bypass} concern that becomes important
for the SAE ablation experiments: in ChooseAttr, both candidate attribute values appear explicitly
in the question string (e.g., ``Is it red or blue?''). Consequently, the model can partially answer
via linguistic priors and object-name associations, without relying on any visual feature.
This dilutes the causal leverage of interventions that remove visual representations.


% ================================================================
\section{Attention Knockout: Reproduction and Extension}
\label{sec:knockout}
% ================================================================

\subsection{Method}

Attention knockout, as introduced by \citet{crossmodal2024}, is a causal probing technique that
zeroes out a selected subset of attention weights during a forward pass to block information flow
between two groups of token positions.
Formally, for a chosen layer $\ell$ and a pair of token sets $(S_{\mathrm{src}},
S_{\mathrm{tgt}})$, the knockout forces $\mathbf{A}^{(\ell)}_{i,j} = 0$ for all $j \in
S_{\mathrm{src}},\ i \in S_{\mathrm{tgt}}$ (pre-softmax row-masking).
Ablation effect is measured as the \emph{margin drop}:
\begin{equation}
  \Delta_{\mathrm{margin}} = \bigl[\log P_{\mathrm{base}}(\text{true}) - \log P_{\mathrm{base}}(\text{false})\bigr]
  - \bigl[\log P_{\mathrm{ko}}(\text{true}) - \log P_{\mathrm{ko}}(\text{false})\bigr].
  \label{eq:margin_drop}
\end{equation}
A positive $\Delta_{\mathrm{margin}}$ indicates that blocking the specified attention path reduces
the model's advantage for the correct answer.

We study two attention-flow directions:
\begin{itemize}[noitemsep]
  \item \textbf{Image$\to$Question}: image patch tokens ($S_{\mathrm{src}}$) attend to question
    tokens ($S_{\mathrm{tgt}}$). This measures how much information the question tokens draw from
    image patches at each layer.
  \item \textbf{Image$\to$Last}: image tokens to the final token (the prediction position).
\end{itemize}

Only items correctly answered by the baseline model are included (\texttt{filter\_correct=True}),
to ensure the metric reflects genuine disruption rather than pre-existing errors.
Effect sizes are reported as Cohen's $d$ from paired $t$-tests across items.

\subsection{Experimental Protocol}

We ran two independent knockout sweeps:

\begin{itemize}[noitemsep]
  \item \textbf{Run 1} (``knockout\_color\_run1''): $n = 510$ correctly-answered color items;
    all 32 layers evaluated.
  \item \textbf{Run 2} (``knockout\_run2\_fixed''): $n = 810$ correctly-answered items
    (all attribute categories); all 32 layers; logprob normalisation fixed.
\end{itemize}

\subsection{Results: Image$\to$Question Flow}

The Image$\to$Question flow shows a clear, reproducible two-peak structure.
Table~\ref{tab:knockout_iq} reports the top layers from both runs.

\begin{table}[ht]
\centering
\caption{Image$\to$Question attention knockout: top layers by margin drop.
  Both runs agree closely on ranking and effect sizes.}
\label{tab:knockout_iq}
\begin{tabular}{lrrrr}
\toprule
Layer & \multicolumn{2}{c}{Run 1 ($n=510$)} & \multicolumn{2}{c}{Run 2 ($n=810$)} \\
\cmidrule(lr){2-3}\cmidrule(lr){4-5}
      & $\Delta_{\mathrm{margin}}$ & $d$ & $\Delta_{\mathrm{margin}}$ & $d$ \\
\midrule
\textbf{0}  & \textbf{0.541} & \textbf{0.833} & \textbf{0.509} & \textbf{0.806} \\
\textbf{11} & \textbf{0.173} & \textbf{0.601} & \textbf{0.169} & \textbf{0.605} \\
8           & 0.076 & 0.435 & 0.064 & 0.414 \\
10          & 0.054 & 0.292 & 0.068 & 0.381 \\
12          & 0.054 & 0.337 & 0.043 & 0.308 \\
14          & 0.033 & 0.324 & 0.032 & 0.361 \\
4           & 0.026 & 0.488 & 0.018 & 0.347 \\
\bottomrule
\end{tabular}
\end{table}

\paragraph{Interpretation.}
Layers 0 and 11 stand apart from all others by a large margin. The gap between layer~0
($d = 0.83$) and layer~11 ($d = 0.60$) is itself substantial, and both substantially exceed the
third-ranked layer ($d \approx 0.44$). This two-peak structure is consistent across both runs
and suggests \emph{two distinct integration phases}:

\begin{enumerate}[noitemsep]
  \item \textbf{Layer 0} — \emph{early cross-modal registration}: the very first transformer layer
    already performs substantial visual grounding, consistent with the hypothesis that attention at
    layer~0 processes raw vision-encoder features before any higher-level language processing.
  \item \textbf{Layer 11} (approximately 1/3 depth) — \emph{higher-level visual-semantic binding}:
    a second wave of cross-modal integration occurs in the early-middle layers, likely corresponding
    to semantic feature composition.
\end{enumerate}

Layer~31 shows $\Delta_{\mathrm{margin}} = 0.0$ by construction (the final layer's output is
already committed before the blocking can propagate).

\subsection{Results: Image$\to$Last Flow}

The Image$\to$Last flow is considerably weaker than Image$\to$Question.
Strikingly, several layers (notably layers 8, 10, 15, 17, and 27) exhibit \emph{negative} margin
drops — blocking those paths \emph{improves} accuracy, suggesting they carry distracting or noisy
information to the final token. The strongest positive effect is at layer~0
($d = 0.20$, Run 2), substantially below the Image$\to$Question signal.
This indicates that question tokens are the primary conduit through which image information is
integrated, rather than a direct path to the last token.

\subsection{Layer Selection for SAE Training}

Based on these results, \textbf{layers 0 and 11} were selected as the two target sites for SAE
training. The attention-output (``\texttt{attn\_out}'') hook was chosen as the activation site,
as it most directly captures the output of the attention mechanism that the knockout probes.


% ================================================================
\section{Sparse Autoencoder Experiments}
\label{sec:sae}
% ================================================================

\subsection{Hypothesis}

The central hypothesis of the SAE phase is:

\begin{quote}
\emph{The causal Image$\to$Question information flow identified at layers 0 and 11 is mediated by
sparse, interpretable features in the residual stream / attention output at those layers.
Ablating the subset of features most selectively associated with visual attribute values should
partially replicate the behavioural disruption seen in the attention knockout, while ablating
matched random features should not.}
\end{quote}

\subsection{SAE Architecture}
\label{sec:sae_arch}

The SAE is a one-hidden-layer autoencoder with ReLU activations:
\begin{equation}
  \mathbf{z} = \mathrm{ReLU}\bigl(\mathbf{W}_{\mathrm{enc}}(\mathbf{x} - \mathbf{b}_{\mathrm{pre}})
    + \mathbf{b}_{\mathrm{enc}}\bigr),
  \qquad
  \hat{\mathbf{x}} = \mathbf{W}_{\mathrm{dec}}\mathbf{z} + \mathbf{b}_{\mathrm{dec}},
  \label{eq:sae}
\end{equation}
trained with a sparsity-penalised MSE loss:
\begin{equation}
  \mathcal{L} = \|\mathbf{x} - \hat{\mathbf{x}}\|_2^2 + \lambda \|\mathbf{z}\|_1.
  \label{eq:sae_loss}
\end{equation}

Through the course of the project the architecture was progressively refined.
Table~\ref{tab:sae_versions} summarises the two major versions.

\begin{table}[ht]
\centering
\caption{SAE architecture versions. V1 was used in runs 1--14; V2 is the current standard.}
\label{tab:sae_versions}
\begin{tabular}{lll}
\toprule
Component & V1 (early experiments) & V2 (current, Mar 2026) \\
\midrule
$n_{\mathrm{features}}$ & 32\,768 & 4\,096 \\
$\lambda$ (\texttt{l1\_coeff}) & $10^{-3}$ & $5 \times 10^{-4}$ \\
Encoder pre-bias $\mathbf{b}_{\mathrm{pre}}$ & absent & present (learnable) \\
Decoder column normalisation & absent & $\|\mathbf{w}_k\|_2 = 1$ after each step \\
Dead-feature tracking & absent & \texttt{dead\_feature\_fraction} logged \\
LR schedule & constant & cosine annealing \\
\bottomrule
\end{tabular}
\end{table}

\subsection{Training Data and Activation Collection}

SAEs were trained on LLaVA activations collected from the ChooseAttr dataset.
For each item, the model was run with the full prompt (image + question) and activations were
extracted at the target layer and activation site.
The activation \emph{position type} determines which sequence positions are collected:

\begin{itemize}[noitemsep]
  \item \texttt{attribute}: tokens corresponding to the attribute-describing region of the question
    (e.g., the colour words ``red or blue'').
  \item \texttt{question}: all question tokens.
  \item \texttt{all}: all positions (image + question + generation).
  \item \texttt{image}: the 576 visual patch token positions (added in March 2026).
\end{itemize}

Training set size for the V1 SAE: approximately 1\,000 ChooseAttr samples yielding
$\sim$6\,000--8\,000 activation vectors at \texttt{attribute} positions.

\subsection{Feature Selection}

Following SAE training, features are ranked by their differential activation between
correctly-answered and incorrectly-answered samples.
Three selection methods were investigated:

\begin{itemize}[noitemsep]
  \item \texttt{ratio}: $\mu_{\mathrm{correct}} / \mu_{\mathrm{incorrect}}$.
  \item \texttt{abs\_diff}: $|\mu_{\mathrm{correct}} - \mu_{\mathrm{incorrect}}|$.
  \item \texttt{causal\_hybrid}: a weighted combination of statistical and causal scores.
\end{itemize}

The top-$k$ ranked features (typically $k = 50$ or $k = 200$) are designated as
\emph{binding features}.

\subsection{Ablation Protocol}

The 3-condition ablation test compares:
\begin{enumerate}[noitemsep]
  \item \textbf{Baseline}: unmodified forward pass.
  \item \textbf{Binding ablation}: the top-$k$ binding features are zeroed in the SAE
    reconstruction.
  \item \textbf{Random control}: a matched set of $k$ randomly-sampled features is zeroed.
\end{enumerate}

Two intervention modes were used:
\begin{itemize}[noitemsep]
  \item \textbf{Residual mode} (legacy):
    $\mathbf{out} = \mathbf{a} + \bigl(\hat{\mathbf{x}}_{\backslash S} - \hat{\mathbf{x}}\bigr)$,
    where $S$ is the ablated feature set.
    This subtracts the selected features' contribution while preserving the reconstruction error.
  \item \textbf{Replace mode} (current standard):
    $\mathbf{out} = \hat{\mathbf{x}}_{\backslash S}$.
    This replaces the full activation with the SAE reconstruction minus the selected features,
    discarding the reconstruction error. This is a harder, more interpretable intervention.
\end{itemize}

All active configurations were standardised to \texttt{replace} mode in March 2026.

\subsection{Complete Experiment Catalogue}

Over the course of the project, 18 SAE ablation experiments were completed or partially run
(Table~\ref{tab:sae_catalogue}).

\begin{longtable}{p{5.5cm}p{2.2cm}p{2.2cm}rr}
\caption{Complete SAE ablation experiment catalogue. $\Delta_{\mathrm{bind}}$ and
  $\Delta_{\mathrm{rand}}$ are the mean margin drops for binding and random conditions respectively.
  ``Null'' indicates no significant separation between conditions.}
\label{tab:sae_catalogue}\\
\toprule
Experiment & Layer / Site & Mode / Position & $\Delta_{\mathrm{bind}}$ & $\Delta_{\mathrm{rand}}$ \\
\midrule
\endfirsthead
\multicolumn{5}{c}{\small\textit{(Table \ref{tab:sae_catalogue} continued)}}\\
\toprule
Experiment & Layer / Site & Mode / Position & $\Delta_{\mathrm{bind}}$ & $\Delta_{\mathrm{rand}}$ \\
\midrule
\endhead
\bottomrule
\endfoot
exp\_run1 & 12 / residual & residual / attr & $-1.71 \times 10^{-5}$ & $-1.71 \times 10^{-5}$ \\
exp\_run2\_all\_residual & 12 / residual & residual / all & $4.40 \times 10^{-4}$ & varies \\
exp\_run3\_attr\_residual\_weak & 12 / residual & residual / attr & $\approx 0$ & $\approx 0$ \\
exp\_run4 & 12 / residual & residual / attr & $2.4 \times 10^{-5}$ & $8.8 \times 10^{-5}$ \\
sae\_q\_layer0 & 0 / residual & residual / attr & $-4.2 \times 10^{-5}$ & $7.7 \times 10^{-5}$ \\
sae\_q\_layer11 & 11 / residual & residual / attr & $3.7 \times 10^{-5}$ & $5.99 \times 10^{-4}$ \\
sweep (residual, $\delta{=}1.0$) & 0 / residual & residual / attr & $1.5 \times 10^{-6}$ & $8.8 \times 10^{-5}$ \\
sweep (residual, $\delta{=}2.0$) & 0 / residual & residual / attr & $-3.8 \times 10^{-6}$ & $4.4 \times 10^{-4}$ \\
sweep (replace, $\delta{=}1.0$) & 0 / residual & replace / attr & $-1.4 \times 10^{-3}$ & $-1.4 \times 10^{-3}$ \\
sweep (replace, $\delta{=}2.0$) & 0 / residual & replace / attr & \multicolumn{2}{c}{\textit{interrupted}} \\
first\_pass\_layer11\_residual & 11 / residual & replace / attr & $1.4 \times 10^{-3}$ & $1.4 \times 10^{-3}$ \\
rerun\_layer11\_attn\_out & 11 / attn\_out & replace / attr & $-2.1 \times 10^{-3}$ & $-2.0 \times 10^{-3}$ \\
layer11\_attn\_out\_replace\_color & 11 / attn\_out & replace / attr & \multicolumn{2}{c}{\textit{null}} \\
layer0\_attn\_out\_replace\_color & 0 / attn\_out & replace / attr & \multicolumn{2}{c}{\textit{null}} \\
layer0\_residual\_replace\_color & 0 / residual & replace / attr & \multicolumn{2}{c}{\textit{null}} \\
layer0\_attn\_out\_diff & 0 / attn\_out & replace / attr & \multicolumn{2}{c}{\textit{null}} \\
layer11\_attn\_out\_replace\_material & 11 / attn\_out & replace / attr & \multicolumn{2}{c}{\textit{null}} \\
with\_holdout\_layer0\_attn\_out & 0 / attn\_out & replace / attr & \multicolumn{2}{c}{\textit{null}} \\
\end{longtable}

\noindent In every completed experiment the binding condition is statistically indistinguishable
from the random control. In the most rigorously controlled run
(``rerun\_layer11\_attn\_out'', $n = 877$ items, 15 matched random-feature sets),
the empirical $p$-values for both accuracy drop and margin drop were 1.0.

\subsection{SAE Reconstruction Quality}

Table~\ref{tab:sae_recon} reports reconstruction metrics from the \texttt{reconstruction\_eval.json}
files of several representative experiments.

\begin{table}[ht]
\centering
\caption{SAE reconstruction quality metrics. All experiments use $n_{\mathrm{features}} = 32\,768$
  and V1 architecture unless noted.}
\label{tab:sae_recon}
\begin{tabular}{lrrrr}
\toprule
Experiment & MSE & Norm. MSE & Expl. Var. & Dead features \\
\midrule
layer0\_attn\_out (with holdout) & $2.15 \times 10^{-6}$ & 0.012 & 98.8\% & 44.2\% \\
layer0\_attn\_out\_diff           & $1.63 \times 10^{-6}$ & 0.009 & 99.1\% & 42.7\% \\
layer0\_residual                  & $4.40 \times 10^{-6}$ & 0.004 & 99.6\% & 0.8\%  \\
layer11\_attn\_out (material)     & $2.27 \times 10^{-4}$ & 0.023 & 97.7\% & 6.1\%  \\
\bottomrule
\end{tabular}
\end{table}

The absolute MSE values are low, and explained variance exceeds 98\% in all cases, which appears
healthy. However, two concerns stand out:
(i) the \texttt{attn\_out} experiments at layer~0 show ${\sim}44\%$ dead features, which is
severely high and indicates that the L1 penalty is over-regularising;
(ii) the mean $L_0$ norm is ${\sim}10\,000$--14\,000 active features per activation vector, far
above the sparse regime targeted by the L1 penalty, suggesting the penalty is simultaneously
too strong (killing features) and insufficient (not producing true sparsity in the survivors).


% ================================================================
\section{Full-Latent Ceiling Analysis}
\label{sec:ceiling}
% ================================================================

After observing a consistent null result across 18 experiments, we designed a diagnostic
\emph{full-latent ceiling ablation} to determine the maximum possible causal effect that the SAE
can produce at each layer and position type.
This experiment (\texttt{10\_full\_latent\_ablation.py}) runs the 3-condition ablation but zeroes
all $n_{\mathrm{features}} = 32\,768$ features simultaneously,
effectively asking: ``if we use the SAE to replace the activation with its reconstruction and then
zero everything, how much does behaviour change?''

\subsection{Ceiling Results}

Table~\ref{tab:ceiling} reports the full-latent ceiling results.

\begin{table}[ht]
\centering
\caption{Full-latent ceiling ablation results ($n = 128$ color items, \texttt{replace} mode).
  Relative perturbation $= \|\hat{\mathbf{x}} - \mathbf{x}\|_2 / \|\mathbf{x}\|_2$.
  The attention knockout margin drops from Table~\ref{tab:knockout_iq} are shown for reference.}
\label{tab:ceiling}
\begin{tabular}{llrrrrr}
\toprule
Layer & Position & Acc.\ drop & Margin drop & Rel.\ pert. & KO margin drop \\
\midrule
11 & \texttt{attribute} & $+3.1\%$ & 0.046 & 0.003 (0.3\%) & 0.17 \\
11 & \texttt{all}       & $+3.1\%$ & 0.169 & 0.999 (99.9\%) & 0.17 \\
0  & \texttt{attribute} & $-0.8\%$ & 0.022 & 0.003 (0.3\%) & 0.54 \\
0  & \texttt{all}       & $-0.8\%$ & 0.314 & 1.016 (101.6\%) & 0.54 \\
\bottomrule
\end{tabular}
\end{table}

\subsection{Interpretation}

The ceiling results are definitive on two points:

\paragraph{Attribute-token positions carry essentially no visual signal.}
At both layers, zeroing all 32\,768 SAE features at \texttt{attribute} token positions produces only
a 0.3\% relative perturbation.
This is not a failure of the SAE to reconstruct well — it reflects the fact that the
\emph{activations at attribute text token positions contain almost no visual information to begin
with}, in the context where the SAE is computing a reconstruction for those tokens.
Feature ablation at question-attribute text tokens is therefore undetectable by design.

\paragraph{The SAE captures the causal signal when ablating all positions.}
When all positions are included, the relative perturbation is $\approx 100\%$ at both layers (the
SAE reconstruction norm matches the original activation norm), and the resulting margin drops are
0.314 (layer~0) and 0.169 (layer~11).
These values are in close correspondence with the attention knockout margin drops of 0.54 and 0.17
respectively: layer~11 matches almost exactly, and layer~0 is within a factor of~2.
This confirms that \emph{the SAE does capture the causal signal} when operating at the correct
positions, and the null result is not a fundamental failure of the SAE to encode visual information.


% ================================================================
\section{Analysis of the Null Result}
\label{sec:null}
% ================================================================

The ceiling analysis, combined with theoretical reasoning, points to six distinct methodological
sources for the null result. We describe each in turn.

\subsection{Position Mismatch (Primary Cause)}

The most fundamental issue is a mismatch between the position at which the attention knockout
operates and the position at which the SAE ablation is applied.

The attention knockout blocks question tokens from attending to image tokens at layer $\ell$.
The causal information being blocked is a quantity computed at the \emph{image token positions} —
specifically, the key-value pairs from image tokens that question tokens would read.
In contrast, SAE ablation in all 18 experiments was applied at \emph{attribute text token
positions} — the target of the attention, not the source.

The ceiling analysis confirms this diagnosis: ablating at \texttt{attribute} positions yields 0.3\%
perturbation because those positions simply do not carry the visual signal that the knockout is
blocking. The correct intervention is to ablate at \textbf{image token positions}.

\subsection{Single-Layer vs.\ Persistent Information Barrier}

Even with correct position targeting, there is a fundamental asymmetry between the two
interventions:
\begin{itemize}[noitemsep]
  \item \textbf{Attention knockout}: zeroes attention weights at layer $\ell$, creating a
    \emph{persistent information barrier} — question tokens cannot read image tokens at layer $\ell$,
    and the resulting deficit propagates through all subsequent layers.
  \item \textbf{SAE ablation}: modifies the \emph{output} of the attention mechanism at layer $\ell$,
    an additive perturbation to the residual stream. At layers $\ell+1$ through 31, the model can
    still attend directly to image tokens, re-reading the visual information that was disrupted.
\end{itemize}
This means that even perfect SAE ablation at image token positions at a single layer cannot fully
replicate the knockout effect; the model compensates by re-reading image tokens at subsequent layers.
The ceiling results corroborate this: the \texttt{all}-position ceiling at layer~0 produces a margin
drop of 0.314 vs.\ the knockout's 0.54 — substantial but not matching — which is precisely what we
would expect if the model partially recovers across layers 1--31.

\subsection{Insufficient SAE Training Data}

The V1 SAE was trained on activations from $\sim$1\,000 ChooseAttr samples, yielding approximately
6\,000--8\,000 activation vectors at attribute token positions.
Modern SAE work~\citep{bricken2023monosemanticity,templeton2024scaling} routinely trains on
millions of activation vectors.
For a 32\,768-feature SAE, the training set of $\sim$7\,000 vectors is orders of magnitude too small:
there are fewer training examples than there are features to learn.
The 44\% dead feature rate observed in the \texttt{attn\_out} experiments is a direct consequence of
this data starvation. The surviving features are likely overloaded and polysemantic, meaning that
features selected as ``colour-discriminating'' may be tracking multiple unrelated concepts.

\subsection{Feature Selection Confound}

The feature selection criterion compares mean activations between correctly and incorrectly answered
items at attribute token positions.
A critical confound is that incorrect answers in the ChooseAttr dataset correlate with
\emph{object size}: 20.6\% of color items contain objects covering less than 1\% of the image area,
and such items are systematically harder to answer correctly.
Features selected by this criterion may therefore be tracking visual difficulty (small object
visibility, image quality, object occlusion) rather than the attribute itself.

\subsection{Language-Side Bypass}

As noted in Section~\ref{sec:dataset}, both answer options appear in the question string in
ChooseAttr.
This allows the model to achieve a non-trivial accuracy without using visual features at all —
via language priors and object-name statistics.
When a portion of the model's performance is language-driven, ablating visual features produces
a correspondingly smaller effect, further diluting any causal signal from feature ablation.

\subsection{SAE Architecture Deficiencies (V1)}

The V1 SAE lacked several components that have been identified as important in subsequent
mechanistic interpretability work:
\begin{itemize}[noitemsep]
  \item No encoder pre-bias $\mathbf{b}_{\mathrm{pre}}$ (subtracted from input before encoding),
    leading to suboptimal feature geometry.
  \item No decoder column normalisation, allowing feature magnitudes to drift and making
    the learned representations less interpretable.
  \item No learning rate schedule; training with a constant LR.
  \item The $\ell_1$ penalty applied as $\lambda \cdot \mathrm{mean}(|\mathbf{z}|)$ rather than a
    per-feature formulation, meaning that increasing $n_{\mathrm{features}}$ reduces effective
    sparsity pressure.
\end{itemize}


% ================================================================
\section{Methodological Improvements (March 2026)}
\label{sec:fixes}
% ================================================================

In response to the ceiling analysis, a series of targeted methodological fixes were implemented
across Plans A--D of a staged improvement roadmap.

\subsection{Plan A: Ceiling Validation and Replace-Mode Standardisation}

The full-latent ceiling experiment (\texttt{10\_full\_latent\_ablation.py}) was designed and
executed, yielding the definitive diagnosis described in Section~\ref{sec:ceiling}.
Concurrently, all nine active experiment configurations were audited and standardised to
\texttt{replace} mode (previously a mix of \texttt{residual} and \texttt{replace} modes existed
across configs, making comparisons unreliable).
A \texttt{grep} for \texttt{"mode: residual"} across all active configs returns zero hits.

\subsection{Plan C: Image-Token Position Type}

A new \texttt{image} position type was implemented in \texttt{ActivationCollector}
(\texttt{sae\_experiments/data/activation\_collector.py}) and \texttt{FeatureAblator}
(\texttt{sae\_experiments/ablation/feature\_ablator.py}).
This position type selects the 576 visual patch token indices starting at the image placeholder
position, computed as:
\begin{equation}
  \text{image\_range} = [\text{img\_placeholder\_idx},\ \text{img\_placeholder\_idx} +
    \text{image\_token\_count})
\end{equation}
consistent with the existing \texttt{get\_image\_token\_range} utility in
\texttt{utils/knockout\_utils.py}.
This allows future experiments to ablate directly at the source of the visual signal rather than
at the textual targets, directly addressing the primary methodological flaw identified in
Section~\ref{sec:null}.
The \texttt{image} position type is set as the default in the V2 configuration.

\subsection{Plan D: SAE Architecture Fixes}

All four architectural deficiencies listed in Section~\ref{sec:null} were fixed in the V2 SAE:
\begin{itemize}[noitemsep]
  \item \textbf{Encoder pre-bias}: a learnable $\mathbf{b}_{\mathrm{pre}}$ parameter was added,
    subtracted from the input before the encoding linear transform.
  \item \textbf{Decoder column normalisation}: \texttt{normalize\_decoder()} is called after every
    optimiser step, constraining each decoder column to unit norm.
  \item \textbf{Dead-feature tracking}: \texttt{dead\_feature\_fraction} is logged per epoch in the
    training history, enabling monitoring and early stopping based on feature utilisation.
  \item \textbf{Cosine annealing LR schedule}: \texttt{CosineAnnealingLR} is applied in
    \texttt{SAETrainer}, replacing the constant learning rate.
\end{itemize}
The V2 configuration (\texttt{configs/sae\_layer0\_attn\_out\_v2.yaml}) reduces $n_{\mathrm{features}}$
from 32\,768 to 4\,096 and $\lambda$ from $10^{-3}$ to $5 \times 10^{-4}$, to better match
the available training data scale.

\subsection{Codebase Refactoring}

Extensive refactoring was performed to consolidate previously duplicated utility functions
into shared modules (Table~\ref{tab:refactor}).

\begin{table}[ht]
\centering
\caption{Centralised utility functions after refactoring (January--March 2026).}
\label{tab:refactor}
\begin{tabular}{lll}
\toprule
Function & Module & Purpose \\
\midrule
\texttt{resolve\_dtype} & \texttt{utils/config\_utils.py} & dtype resolution (was duplicated in 5 scripts) \\
\texttt{get\_target\_module} & \texttt{utils/hook\_utils.py} & layer + site to nn.Module lookup \\
\texttt{estimate\_image\_token\_count} & \texttt{utils/knockout\_utils.py} & dynamic image token count \\
\texttt{get\_question\_token\_range} & \texttt{utils/knockout\_utils.py} & canonical sublist-search \\
\texttt{sequence\_logprob} & \texttt{utils/knockout\_utils.py} & log-prob scoring utility \\
\texttt{setup\_experiment} & \texttt{utils/script\_utils.py} & experiment boilerplate (7 scripts) \\
\texttt{load\_llava\_components} & \texttt{utils/script\_utils.py} & model loading boilerplate \\
\texttt{load\_sae} & \texttt{utils/script\_utils.py} & SAE checkpoint loading \\
\bottomrule
\end{tabular}
\end{table}

Dead code — including unused methods in \texttt{AblationExperiment} and \texttt{HypothesisTester}
— was removed, reducing the total codebase complexity.

% ================================================================
\section{PaliGemma-2 Extension}
\label{sec:paligemma}
% ================================================================

A new branch (\texttt{paligemma-2}) was created to extend the experimental pipeline to a second
MLLM: \textbf{Google PaliGemma-2-3B} (\texttt{google/paligemma2-3b-pt-224}),
a 3-billion-parameter model with a SigLIP vision encoder and Gemma-2 language backbone
($d_{\mathrm{model}} = 2304$, $256$ image tokens at 224\,px resolution).

The extension involves three main additions:

\paragraph{PaliGemma dataset adapter.}
\texttt{sae\_experiments/data/paligemma\_dataset.py} implements an adapter that reformats
ChooseAttr items into PaliGemma's prompt convention (SFT prefix and suffix tokens).

\paragraph{PaliGemma hook utilities.}
\texttt{sae\_experiments/utils/paligemma\_hooks.py} provides model-specific layer targeting,
handling differences in the attention module naming convention between LLaMA-2 and Gemma-2.

\paragraph{GemmaScope SAE integration.}
\texttt{sae\_experiments/models/gemma\_scope\_sae.py} implements a JumpReLU SAE
weight-compatible with the publicly released \textbf{Gemma Scope} checkpoints~\citep{gemmascope2024}:
\begin{equation}
  \mathbf{z} = \mathrm{ReLU}(\mathbf{x}_{\mathrm{cent}} \mathbf{W}_{\mathrm{enc}} + \mathbf{b}_{\mathrm{enc}})
    \cdot \mathbf{1}[\mathbf{x}_{\mathrm{cent}} \mathbf{W}_{\mathrm{enc}} + \mathbf{b}_{\mathrm{enc}} > \boldsymbol{\theta}],
\end{equation}
where $\boldsymbol{\theta}$ is a learnable per-feature JumpReLU threshold, and
$\mathbf{x}_{\mathrm{cent}} = \mathbf{x} - \mathbf{b}_{\mathrm{dec}}$.
This enables loading pre-trained Gemma Scope weights directly, bypassing the need to train an SAE
from scratch and circumventing the training-data-scale problem that plagued the LLaVA experiments.

\paragraph{Significance.}
The PaliGemma-2 extension represents a strategic pivot: rather than training a custom SAE on a
small ChooseAttr corpus, the plan is to leverage the publicly available, large-scale Gemma Scope
SAEs (trained on billions of tokens) and apply the same knockout + feature ablation pipeline.
This directly addresses two of the six null-result causes simultaneously:
insufficient training data (pre-trained SAE) and polysemantic features (Gemma Scope SAEs achieve
modern sparsity standards).


% ================================================================
\section{Proposed Alternative Task Format}
\label{sec:task}
% ================================================================

As identified in Section~\ref{sec:null}, the ChooseAttr forced-choice format introduces a
language-side bypass.
We evaluated three alternative task formats:

\begin{enumerate}
  \item \textbf{Open-ended logprob scoring} (recommended near-term):
    Score the log-probability of the single correct answer token given the open-ended prompt
    \texttt{"<image>\textbackslash nWhat color is the \{obj\}? Answer with one word."}.
    This eliminates the bypass by removing the false-option text from the prompt.
    The \texttt{sequence\_logprob} utility in \texttt{utils/knockout\_utils.py} supports this
    directly. Of the 602 color rows, 599/602 have single-word answers; only 3 rows involve
    multi-word answers (\emph{dark brown}, \emph{dark blue}) and can be excluded without
    meaningful loss.
    This format requires approximately 30 lines of new code and no re-annotation.

  \item \textbf{GQA QueryAttr} (medium effort):
    The file \texttt{datasets/GQA\_val\_correct\_question\_with\_positionQuery\_QueryAttr.csv}
    exists in the repository but contains only \texttt{positionQuery} questions; authentic
    open-ended queryAttr questions must be fetched separately from HuggingFace.

  \item \textbf{Image-token activation probing} (medium-high effort):
    Collect SAE feature activations at image token positions for images depicting the target
    attribute, with no language prompt. This tests whether features fire for visual content
    independently of language context and is theoretically the strongest test.
    It directly replicates the knockout by removing the attribute signal at its source.
    Supported by \citet{tong2024iclr} and the SAE-V approach~\citep{saev2025}.
\end{enumerate}

\noindent We recommend implementing format (1) as the immediate next step, as it eliminates the
language bypass while reusing the existing dataset, evaluation infrastructure, and statistical
pipeline.


% ================================================================
\section{Summary of Quantitative Results}
\label{sec:summary}
% ================================================================

\subsection{What We Know With High Confidence}

\begin{enumerate}
  \item \textbf{Layers 0 and 11 are the dominant causal pathways} for Image$\to$Question
    information flow in LLaVA-v1.5-7b on ChooseAttr, with Cohen's $d = 0.83$ and $0.60$
    respectively. This finding replicates across two independent knockout sweeps with different
    sample sizes ($n = 510$ and $n = 810$).

  \item \textbf{The SAE captures the causal signal when operating at all positions}.
    The full-latent ceiling at \texttt{all} positions produces margin drops of 0.314 (layer~0)
    and 0.169 (layer~11), closely matching the knockout values of 0.54 and 0.17 respectively
    (particularly for layer~11, where the match is near-exact).

  \item \textbf{Attribute text token positions carry essentially no visual signal} (0.3\%
    relative perturbation in the full-latent ceiling). This is a property of the information
    routing in the model, not a deficiency of the SAE.

  \item \textbf{The null result in 18 SAE ablation experiments is methodological}, arising
    primarily from the position mismatch between the attention knockout (which operates at
    the source of visual information) and the SAE ablation (which operated at target
    text-token positions).
\end{enumerate}

\subsection{What Remains Open}

\begin{enumerate}
  \item Whether SAE features at \emph{image token positions} at layers 0 and 11 show specific
    activations for visual attributes (color, material, etc.).
  \item Whether ablating such features at image token positions produces a causal effect on
    attribute-binding behaviour.
  \item Whether the two-peak structure in Image$\to$Question flow corresponds to qualitatively
    different types of feature (raw colour encoding at layer~0 vs.\ semantic attribute binding
    at layer~11).
  \item Whether the findings replicate in PaliGemma-2 with Gemma Scope pre-trained SAEs,
    providing a model-general picture of cross-modal information routing.
\end{enumerate}


% ================================================================
\section{Prioritised Next Steps}
\label{sec:next}
% ================================================================

The following steps are ordered by expected impact and implementation cost:

\begin{enumerate}[label=\textbf{E\arabic*.}]

  \item \textbf{Knockout-guided feature identification at image token positions (Plan E).}
    Run paired forward passes with and without the layer-0 attention knockout, collecting
    SAE feature activations at image token positions under both conditions.
    The difference in feature activation patterns between knockout and non-knockout conditions
    directly identifies which features carry the Image$\to$Question information.
    This bypasses the feature-selection confound entirely and provides a direct supervision
    signal from the causal intervention.

  \item \textbf{Validate image-position ablation with the V2 SAE.}
    Using the \texttt{image} position type (Plan C) and the V2 SAE architecture (Plan D),
    run the 3-condition ablation experiment at image token positions.
    Compare the margin drop against the full-latent ceiling to quantify what fraction of the
    causal signal is captured by the ranked binding features.

  \item \textbf{Switch to open-ended logprob scoring (Plan B / Section~\ref{sec:task}).}
    Implement format (1) from Section~\ref{sec:task} to eliminate the language-side bypass.
    This can be done as a parallel evaluation alongside ChooseAttr to isolate the contribution
    of the bypass to the null result.

  \item \textbf{Scale SAE training data.}
    Collect activations from a diverse corpus (LLaVA-Instruct or full GQA validation,
    $\geq 100\,000$ activation vectors) to reduce the dead-feature rate and improve feature
    monosemanticity. This is a prerequisite for interpretable feature-level analysis.

  \item \textbf{Deploy PaliGemma-2 + Gemma Scope pipeline (Plan F).}
    Complete the PaliGemma-2 port (branch \texttt{paligemma-2}): run the attention knockout
    sweep to identify causal layers, then load the corresponding Gemma Scope SAE checkpoints
    and run the image-position ablation experiment.
    This provides a second data point on whether the two-peak structure generalises across
    architectures.

\end{enumerate}


% ================================================================
\section{Conclusion}
\label{sec:conclusion}
% ================================================================

This project has made substantial progress on understanding cross-modal information flow
in LLaVA-v1.5-7b. The attention knockout reproduction confirms and extends the original
findings~\citep{crossmodal2024}, establishing layers 0 and 11 as the dominant causal
pathways for Image$\to$Question information transfer, with large and reproducible effect sizes.

The SAE ablation phase produced a consistent null result across 18 experiments, but this
null result has been rigorously diagnosed through a ceiling analysis and theoretical reasoning.
The primary cause is a position mismatch: visual information flows through \emph{image token}
positions, but all 18 ablation experiments targeted \emph{attribute text token} positions.
The full-latent ceiling confirms that the SAE captures the causal signal when operating at
all positions, and that the margin drops produced match the knockout magnitudes closely at
layer~11.

A complete set of methodological fixes has been implemented (Plans A--D), including a new
\texttt{image} position type for activation collection and ablation, V2 SAE architecture with
modern best practices, standardised \texttt{replace} mode across all configurations, and
extensive codebase refactoring.
The groundwork for a strategically superior experiment — image-position ablation with
knockout-guided feature identification using pre-trained Gemma Scope SAEs on PaliGemma-2 —
is in place.

The core scientific question remains open and tractable: \emph{which sparse, interpretable
features at the image token positions of early-to-middle transformer layers causally mediate
visual attribute binding in MLLMs?}
We are now in a position to answer it.


% ================================================================
\bibliography{references}
\bibliographystyle{plainnat}
% ================================================================

\begin{thebibliography}{99}

\bibitem[Bricken et al.(2023)]{bricken2023monosemanticity}
Bricken, T., Templeton, A., Batson, J., Chen, B., Jermyn, A., Conerly, T., \ldots\ \& Olah, C.\
  (2023).
Towards monosemanticity: Decomposing language models with dictionary learning.
\textit{Transformer Circuits Thread}.

\bibitem[Crossmodal(2024)]{crossmodal2024}
Anonymous.\ (2024).
Cross-modal information flow in multimodal large language models.
\textit{Preprint / NeurIPS Workshop}.

\bibitem[Gemma Scope(2024)]{gemmascope2024}
Google DeepMind.\ (2024).
Gemma Scope: Open sparse autoencoders everywhere all at once on Gemma 2.
\textit{Preprint}.

\bibitem[Liu et al.(2023)]{liu2023llava}
Liu, H., Li, C., Wu, Q., \& Lee, Y.\ J.\ (2023).
Visual instruction tuning.
\textit{NeurIPS 2023}.

\bibitem[SAE-V(2025)]{saev2025}
Anonymous.\ (2025).
SAE-V: Interpreting vision-language models with sparse autoencoders.
\textit{ICML 2025}.

\bibitem[Templeton et al.(2024)]{templeton2024scaling}
Templeton, A., Conerly, T., Marcus, J., Lindsey, J., Bricken, T., Chen, B., \ldots\ \& Olah, C.\
  (2024).
Scaling and evaluating sparse autoencoders.
\textit{Anthropic Research}.

\bibitem[Tong et al.(2024)]{tong2024iclr}
Tong, S., Liu, Z., Zhai, Y., Ma, Y., LeCun, Y., \& Xie, S.\ (2024).
Eyes wide shut? Exploring the visual shortcomings of multimodal LLMs.
\textit{CVPR 2024}.

\end{thebibliography}

\end{document}
